\section{A proof-producing interface to iPRA}
\label{logic:section}

The repository already supplies a suitable foundational endpoint
\cite{repoCore}.  In
\path{ConciseIPRA/Lean/ConciseIPRA/PR.lean}, the inductive family
\texttt{PR k} represents \emph{definitions} of primitive recursive functions
of arity $k$, not arbitrary total functions.  Its constructors are zero,
successor, projection, composition, and primitive recursion.  The recursion
argument comes first:
\[
 R_{g,h}(0,\vec x)=g(\vec x),\qquad
 R_{g,h}(n+1,\vec x)=h(n,R_{g,h}(n,\vec x),\vec x).
\]
Terms \texttt{PTm n} are applications of these symbols to terms and variables.
The theory \texttt{PRAAx} has their defining equations, successor separation
and injectivity, and induction for quantifier-free formulas.  The surrounding
logic is the intuitionistic natural deduction of \texttt{Deduction.lean}, with
equality, the propositional connectives, and first-order quantifiers.
An actual formalization target is therefore
\[
   \texttt{Prf PRAAx}\;\Gamma\;\varphi,
   \qquad\text{abbreviated }\Gamma\vdash_{\mathrm{iPRA}}\varphi.
\]
This is a statement about a derivation in the object theory.  A Lean proof of
the standard numerical interpretation of $\varphi$ has a different type.
The distinction is essential: the second proof may use induction or choice
that the first theory does not possess.

\subsection{Finite presentations and bounded formulas}

A finite mathematical sort is presented by natural-number codes, a primitive
recursive domain test $D$, and a primitive recursive equality test $E$.
Equality need not be literal equality of codes.  For example, the rational
pairs $(a,b)$ with $b>0$ may use cross multiplication; a proof that the
operations respect this relation replaces any appeal to a quotient type.
Each operation is a \emph{fixed} primitive recursive function, and its domain
preservation and congruence laws are proved in iPRA.  Quantifiers are translated
by
\[
 (\forall a\in A\,\theta)^*=
 \forall u\,(D(u)\longrightarrow\theta^*(u)),\qquad
 (\exists a\in A\,\theta)^*=
 \exists u\,(D(u)\wedge\theta^*(u)).
\]
There is no new sort in the object theory.  A function with extra side
conditions, such as rational inversion, is made total by a default value;
its useful equations are asserted only under the advertised side conditions.

We call a formula \emph{bounded} when all its quantifiers range below specified
primitive recursive terms (each bound depends only on variables outside its
own binder) and all its atomic predicates have primitive
recursive characteristic functions with proved correctness.  This definition
is syntactic.  It does not include a predicate merely because someone knows
that its truth can be decided.

\begin{theorem}[Booleanization of a fixed bounded formula]
\label{logic:booleanization}
For every fixed bounded formula $B(\vec x)$ constructed as above, there is a
primitive recursive Boolean function $b_B(\vec x)$ and an intuitionistic iPRA
derivation
\[
 B(\vec x)\ \longleftrightarrow\ b_B(\vec x)=1.
\]
The construction is effective on a finite syntax tree.  Each resulting
function symbol is fixed before it is used in an object-theory theorem.
\end{theorem}

\begin{proof}
The atomic case is part of the presentation.  Use Boolean conjunction and
disjunction for $\wedge$ and $\vee$.  An implication is represented by
$(1-b_A)\vee b_C$, and falsity by $0$.  The equivalence for implication is
intuitionistically valid because $b_A$ is a Boolean: in the case $b_A=0$ the
antecedent is refutable; in the case $b_A=1$ the implication supplies the
consequent.  These are cases on a decidable numerical value, not an application
of excluded middle to an arbitrary proposition.

For a bounded universal quantifier use the primitive recursive fold
\[
 a(0,\vec x)=1,\qquad
 a(t+1,\vec x)=a(t,\vec x)\wedge b_B(t,\vec x).
\]
Its elementary Boolean laws and the assertion that it is Boolean are proved
by quantifier-free induction.  To avoid concealing a stronger induction in
the explanation, fix the final bound $N$ as a parameter.  Under the hypothesis
$\forall i<N\,b_B(i,\vec x)=1$, induct on the quantifier-free assertion
\[
       t\le N\ \longrightarrow\ a(t,\vec x)=1.
\]
At the step with $t+1\le N$, instantiate the fixed hypothesis at $t$.  The
other case is vacuous.  This proves the forward direction at $t=N$.
For the converse, fix also an index $i$ and prove by quantifier-free induction
on $t$ that
\[
 i<t\ \wedge\ a(t,\vec x)=1
       \ \longrightarrow\ b_B(i,\vec x)=1.
\]
At the successor step, decide whether $i=t$ or $i<t$ and use the defining
equation of the fold.  Universal introduction over $i$ is a logical rule and
is not an additional induction principle.

For a bounded existential quantifier use the success flag of the bounded
search described below.  The search invariant can be made quantifier-free
\emph{directly}: write its assertions about failing prefixes as equations for
the conjunction fold of $1-b_B$.  Thus this step does not presuppose the
existential case being proved.  If the flag accepts, search returns a term
below the bound whose atomic test accepts; existential introduction supplies
the witness.  Conversely, a bounded witness contradicts the failing-prefix
equation if the flag is zero, so the Boolean flag must be one.  Applying the
induction hypothesis for the subformula completes this case.
The metatheoretic structural recursion over the \emph{fixed} formula builds a
finite derivation.  It does not introduce an object-level truth predicate.
\end{proof}

Thus a bounded universal over a coded vector, a finite sum inequality, a
finite matrix identity, or a bounded minimum condition can be used as an
induction invariant after replacement by its Boolean equation.  This argument
does not justify induction on an unbounded formula such as
$\forall m\,R(n,m)$, or on the bare statement that a real-coded object has a
property at every precision.

\subsection{Bounded search supplies its own witness}

Let $p(i,\vec x)$ be a fixed primitive recursive Boolean test.  Search scans
$0,\ldots,N-1$ and returns the first successful index, or a distinguished
failure value.  An implementation may use the pair $(\mathrm{found},j)$,
so no undecidable convention is involved.  At stage $t$ the invariant is
\[
\begin{split}
 &\mathrm{found}=0\quad\Longrightarrow\quad
      \forall i<t\,p(i,\vec x)=0,\\
 &\mathrm{found}=1\quad\Longrightarrow\quad
      j<t\ \wedge\ p(j,\vec x)=1\ \wedge\
      \forall i<j\,p(i,\vec x)=0.
\end{split}
\]
For a direct quantifier-free version, let $a(t)$ be the conjunction fold of
$1-p(i,\vec x)$ for $i<t$, and replace the two universal prefix assertions by
$a(t)=1$ and $a(j)=1$, respectively.  This construction uses only Boolean
primitive recursion, so there is no circular dependence on the existential
part of Theorem~\ref{logic:booleanization}.  Initially the first alternative
holds vacuously.  A stored successful index is retained.  If no index is
stored, the decidable test at $t$ either stores $t$ or extends the failing
prefix.  These are all the cases, so quantifier-free induction proves the
invariant.  The already proved universal-fold equivalence gives its displayed
readable form.
In particular,
\[
 \exists i<N\,p(i,\vec x)=1
   \quad\Longrightarrow\quad
 p(\mathrm{search}(N,\vec x),\vec x)=1
 \ \wedge\ \mathrm{search}(N,\vec x)<N,
\]
where the displayed selector uses a fixed default on failure.
Existence of a bounded witness is enough; unbounded minimization is never
used.  If no bound $N$ is available, this lemma does not provide one.

The same proof gives a useful finite-observation fact.  If two Boolean tables
agree below $N$, their searches below $N$ are equal.  At stage $t<N$ the
algorithms have equal stored states and read the same table entry; induction
therefore preserves equality.  This makes the dependency on finite data
explicit even though the search may stop at an input-dependent position.

\subsection{A bounded-loop compilation theorem}

Fix primitive recursive functions $s_0(\vec x)$, $F(t,s,\vec x)$ and
$T(\vec x)$.  A tuple of state variables is coded by one number.  Define
\[
 s(0,\vec x)=s_0(\vec x),\qquad
 s(t+1,\vec x)=F(t,s(t,\vec x),\vec x),\qquad
 \operatorname{out}(\vec x)=s(T(\vec x),\vec x).
\]
All three are fixed primitive recursive functions.  The following theorem is
the central reusable rule; it states exactly what an invariant proof must
supply.

\begin{theorem}[Bounded-loop rule in iPRA]
\label{logic:boundedloop}
Let $I(t,s,\vec x)$ be a fixed bounded formula and let $\Gamma(\vec x)$ be a
finite context independent of $t$.  Suppose iPRA derives
\[
\begin{gathered}
 \Gamma\vdash I(0,s_0(\vec x),\vec x),\\
 \Gamma,\ t<T(\vec x),\ I(t,u,\vec x)
       \vdash I(t+1,F(t,u,\vec x),\vec x).
\end{gathered}
\]
Then iPRA derives
\[
             \Gamma\vdash
             I(T(\vec x),\operatorname{out}(\vec x),\vec x).
\]
The context $\Gamma$ may contain unbounded assumptions; those assumptions
remain explicit and are not asserted to be decidable.
\end{theorem}

\begin{proof}
Replace $I$ by its proved Boolean equivalent $b_I=1$.  Induct on the
quantifier-free formula
\[
 J(t,\vec x):\quad
 t\le T(\vec x)\ \longrightarrow\ b_I(t,s(t,\vec x),\vec x)=1.
\]
The base follows from the initial contract and the equation for $s(0)$.
For the successor step, assume $J(t,\vec x)$ and $t+1\le T(\vec x)$.
Elementary numerical arithmetic gives $t<T(\vec x)$ and $t\le T(\vec x)$.
The induction hypothesis gives $I(t,s(t,\vec x),\vec x)$.  Substitute the
term $s(t,\vec x)$ for $u$ in the supplied transition derivation.  Its
conclusion, the defining equation for $s(t+1)$, and Booleanization give
$J(t+1,\vec x)$.  This argument is intuitionistic implication introduction;
it needs no assertion about the state after the advertised bound.

The repository's rule \texttt{PRA.qfInd} now gives $\forall t\,J(t,\vec x)$.
Instantiate it at the term $T(\vec x)$ and use reflexivity of $\le$.
Finally use the reverse direction of Booleanization.  Every induction used
in this proof is on a quantifier-free formula.  Substitution, weakening and
cut are already derived operations in \texttt{Deduction.lean}.
\end{proof}

An early exit is handled by a Boolean flag in the state and a transition that
retains the state once the flag is set.  Nested loops are handled from the
inside outward.  The nesting depth is fixed by the program.  No small
complexity bound is needed: a tower of exponentials of primitive recursive
height is permissible when supplied by a fixed primitive recursive
definition and a valid invariant proof.

\begin{theorem}[Finite contract compilation]
\label{logic:compilation}
Consider a fixed finite program built from fixed primitive recursive
assignments, finite tuples, decidable branches, sequential composition and
loops with primitive recursive bounds.  Suppose its leaf lemmas have iPRA
derivations and every loop has the initial and transition derivations of
Theorem~\ref{logic:boundedloop}.  Then its precondition-to-postcondition
derivation can be assembled in iPRA.  If the result is advertised as an
existential witness, the witness is the program's fixed primitive recursive
output term.
\end{theorem}

\begin{proof}
Proceed through the finite construction of the program.  At an assignment,
substitute its defining term into the postcondition and use its supplied
defining equations.  At a sequential composition, substitute the first
output term into the second derivation and cut with the first contract.  At
a branch, the Boolean condition is either $0$ or $1$; the two supplied
branch derivations yield the selected postcondition by intuitionistic
disjunction elimination.  At a loop apply Theorem~\ref{logic:boundedloop}.
These constructions preserve the advertised context and produce a derivation
of the postcondition with the output term substituted for the output
variable.  The usual existential introduction rule then introduces that
term.  This construction takes a finite supplied derivation description;
it is not an assertion that arbitrary valid invariants can be found.
\end{proof}

This theorem deliberately proves a restricted, transparent compilation
result.  It does \emph{not} rely on a general theorem extracting primitive
recursive witnesses from every possible first-order iPRA proof.  Nor does it
infer that a primitive recursive function's true specification is provable
in iPRA.  The supplied derivations are indispensable.

\subsection{Static definitions, clocks and the uniformity boundary}

The phrase ``a primitive recursive code for a real'' can conceal an invalid
universal evaluator.  The following elementary obstruction fixes the issue.

\begin{theorem}[No primitive recursive universal evaluator]
\label{logic:noeval}
There is no primitive recursive $U(e,x)$ whose sections $U(e,\cdot)$ include
every unary primitive recursive function.
\end{theorem}
\begin{proof}
If such $U$ existed, then $d(n)=U(n,n)+1$ would be primitive recursive.
Choose an index $e$ for $d$.  Substitution gives
$d(e)=U(e,e)+1=d(e)+1$, a contradiction in elementary numerical arithmetic.
\end{proof}

There are three sound interfaces.  First, a \emph{static schema} fixes the
primitive recursive approximation functions before stating its iPRA theorem.
Numerical parameters may still vary uniformly.  Second, a theorem may take a
\emph{finite observation table} as an ordinary numerical input, prove the
result from its explicit bounded consistency conditions, and use no evaluator
for an infinite object.  Third, a theorem may take a program code and an
\emph{execution clock} or finite execution trace.  Simulating a fixed
universal machine for a supplied number of steps and checking a finite trace
are primitive recursive.  The theorem says what follows if the simulation
finishes and the trace passes its checks; it does not remove the clock.

For each \emph{fixed} primitive recursive definition $f$, a primitive
recursive trace producer can be built by recursion on that fixed definition.
For zero, successor and projection the trace is immediate.  For composition,
produce the finitely many argument traces, then the outer trace, and
concatenate.  For primitive recursion, primitive-recursively accumulate the
current value together with the trace; append the fixed step program's trace
at each iteration.  Concatenation and sequence coding are primitive recursive,
so the new trace producer is primitive recursive.  Its correctness is proved
using the invariant ``this accumulated trace is valid and ends with this
current value'', a bounded trace predicate.  This is a metatheoretic
construction for each fixed syntax tree, with an iPRA proof for the resulting
program.  Theorem~\ref{logic:noeval} forbids turning it into a single
primitive recursive unclocked evaluator over all such trees.

An ordinary Lean structural evaluator can therefore be total while failing
to be primitive recursive uniformly in its source code.  The same caution
applies to full normalization in a compact proof checker.  For internal
arithmetization use the raw defining equations, or explicit reduction traces,
unless a separate bound on normalization has been proved.  This observation
does not question the repository's interderivability theorems; it prevents
an extra computational conclusion from being silently added to them.

\subsection{Verifiers, producers and actual object proofs}

A verifier with a soundness contract has the form
\[
 V(\vec x,c)=1\quad\Longrightarrow\quad R(\vec x).
\]
It does not by itself prove $R(\vec x)$.  For that, the proposed construction
must additionally give a primitive recursive producer $c=C(\vec x)$ and an
iPRA derivation of $V(\vec x,C(\vec x))=1$ under its advertised hypotheses.
Cut then proves $R$.  A bounded-height point enumerator, for example, can
verify every point it returns without proving that no point lies outside its
search range.  That missing exhaustion argument is a different contract.

Similarly, a Lean theorem saying that an accepted proof certificate yields
\texttt{Prf PRAAx Gamma phi} is a metatheorem about derivability.  It is not an
internal iPRA truth predicate or a global reflection principle.  The existing
\texttt{pra\_iff\_native} transfers a supplied iPRA derivation to the native
calculus, and \texttt{native\_iff\_pra} transfers it back.  Neither theorem
turns a semantic Lean proof into an object-theory derivation.

\subsection{Lean realization and its remaining obligations}

The delivered \texttt{FiniteContracts.lean} proves, in ordinary Lean, the
bounded-loop invariant, least bounded-search correctness, witness recovery
under bounded existence, and dependence only on the finite tested prefix.
The definitions use only \texttt{Std}; they have been compiled with Lean
4.32.0.  These semantic lemmas help test the algorithms and their interfaces,
but the file explicitly does not identify arbitrary Lean predicates with
primitive recursive ones.

The delivered \texttt{ObjectKernel.lean} concatenates the pinned repository's
\texttt{Syntax.lean}, \texttt{Deduction.lean} and \texttt{PR.lean}, retaining
their declarations and proofs.  The imports are replaced by \texttt{Std},
with the notation $\mathbb N$ restored.  This standalone kernel compiles.
\texttt{ObjectBridge.lean} then defines a contract by the data
\[
\begin{array}{l}
 \texttt{pre : Fm PTm n},\quad \texttt{program : PR n},\\
 \texttt{post : Fm PTm (n+1)},\\
 \texttt{certified : [pre] |- post[program(x)/output]}.
\end{array}
\]
It derives the existential output statement and records the exact
\texttt{PRA.qfInd} interface.  These are genuine object-theory derivations.
The proof fields are obligations, not axioms asserting that every program
satisfies its desired contract.

\texttt{ObjectExamples.lean} supplies an explicit primitive recursive symbol
for addition and proves the closed object-iPRA theorem
$\forall x\,(x+0=x)$ by quantifier-free induction, as well as a simple witness
contract.  Thus the scoping, symbol equations, induction rule and witness
introduction have been tested together.  The axiom audit lists no
\texttt{sorryAx} and no \texttt{Classical.choice}; the object kernel uses the
ordinary Lean principles \texttt{propext} and \texttt{Quot.sound} already used
by the upstream development.  This concerns Lean's metatheory and does not
add excluded middle to the embedded derivation rules.

For the full formalization the next modules are precise: construct sequence
and signed-rational symbols; build Booleanization by induction on bounded
formula syntax; prove each loop's finite transition lemmas as
\texttt{Prf PRAAx}; assemble the induction certificates; and instantiate the
contract records with those proofs.  Finally the separately supplied
\texttt{RepositoryBridge.lean} invokes \texttt{pra\_iff\_native}.  That small
adapter has been checked against the source interface but was not compiled
against the entire repository dependency closure.  A complete general
Booleanization compiler and the object-iPRA derivations of the later analytic
applications remain formalization work; they are not hidden in the compiled
examples.
